\begin{figure}[ht]
\centering
\begin{tikzpicture}[x=1cm, y=1cm, font=\small]

% Left panel: the fitted exponential-decay curve through scattered data.
\begin{scope}
  \draw[->, thick] (-0.2, 0) -- (5.6, 0)
    node[right, font=\scriptsize] {$t\ [\si{\hour}]$};
  \draw[->, thick] (0, -0.2) -- (0, 3.6)
    node[above, font=\scriptsize] {$y$};

  % Model y = A exp(-k t) + C with A=2.6, k=0.55, C=0.4 (axis-scaled).
  % Draw the smooth fitted curve.
  \draw[engblue, very thick, samples=80, domain=0:5.2, smooth]
    plot (\x, {0.4 + 2.6*exp(-0.55*\x)});

  % Scattered data points (noisy samples around the curve).
  \fill[engblack] (0.2, 3.05) circle (1.3pt);
  \fill[engblack] (0.6, 2.50) circle (1.3pt);
  \fill[engblack] (1.0, 2.18) circle (1.3pt);
  \fill[engblack] (1.5, 1.70) circle (1.3pt);
  \fill[engblack] (2.0, 1.32) circle (1.3pt);
  \fill[engblack] (2.6, 1.08) circle (1.3pt);
  \fill[engblack] (3.2, 0.86) circle (1.3pt);
  \fill[engblack] (3.9, 0.66) circle (1.3pt);
  \fill[engblack] (4.6, 0.58) circle (1.3pt);
  \fill[engblack] (5.1, 0.50) circle (1.3pt);

  \node[engblue, anchor=west, font=\scriptsize] at (2.7, 1.95)
    {$\hat{y}(t) = A\,e^{-kt} + C$};
  \node[anchor=north, font=\scriptsize] at (2.6, -0.55) {(a) fitted decay curve};
\end{scope}

% Right panel: residual-norm convergence, Gauss-Newton vs gradient descent.
\begin{scope}[xshift=7.5cm]
  \draw[->, thick] (-0.2, 0) -- (4.6, 0)
    node[right, font=\scriptsize] {iteration $k$};
  \draw[->, thick] (0, -0.2) -- (0, 3.6)
    node[above, font=\scriptsize] {$\log_{10}\norm{\vect{r}_{k}}$};

  % Gauss-Newton / LM: fast (near-quadratic) drop.
  \draw[engred, very thick]
    (0, 3.3) -- (0.7, 2.4) -- (1.4, 1.2) -- (2.1, 0.35) -- (2.8, 0.10)
    -- (3.6, 0.06);
  \foreach \p in {(0,3.3),(0.7,2.4),(1.4,1.2),(2.1,0.35),(2.8,0.10),(3.6,0.06)}
    \fill[engred] \p circle (1.4pt);
  \node[engred, anchor=south west, font=\scriptsize] at (2.2, 0.45)
    {Levenberg-Marquardt};

  % Gradient descent: slow geometric drop.
  \draw[enggray, very thick]
    (0, 3.3) -- (0.7, 3.05) -- (1.4, 2.82) -- (2.1, 2.60) -- (2.8, 2.40)
    -- (3.6, 2.18) -- (4.3, 1.98);
  \node[enggray, anchor=south west, font=\scriptsize] at (2.0, 2.55)
    {gradient descent};

  \node[anchor=north, font=\scriptsize] at (2.2, -0.55)
    {(b) residual norm vs iteration};
\end{scope}

\end{tikzpicture}
\caption[Gauss-Newton calibration of a decay model]{Nonlinear
least-squares calibration of the exponential-decay model $\hat{y}(t)
= A\,e^{-kt} + C$. Panel (a): the fitted curve threads the scattered
measurements, each datum contributing a residual $r_{i} = y_{i} -
\hat{y}(t_{i})$ that the calibration drives down in the sum-of-squares
sense. Panel (b): the residual norm against iteration on a log scale.
Levenberg-Marquardt, which uses the Gauss-Newton approximation
$J^{\transpose}J$ to the Hessian, reaches the optimum in a handful of
iterations with the near-quadratic tail of a second-order method;
plain gradient descent on the same objective crawls down the
ill-conditioned valley at the geometric rate set by the condition
number of $J^{\transpose}J$. The contrast is the practical reason
every curve-fitting library defaults to Levenberg-Marquardt rather
than to a first-order method.}
\label{fig:vol02:ch15:gauss-newton-fit}
\end{figure}
